% macros.tex -- notation and theorem environments of the article and its Supplementary Material.
% Loaded by main.tex and supplement.tex after the packages.  One symbol per quantity; symbols that are used in
% a single section of the Supplementary Material only are defined at the top of that section.
%
% =====================================================================================
% NOTATION TABLE (single source of truth)
% -------------------------------------------------------------------------------------
% CONVENTIONS
%   * Sites are ZERO-BASED: \Om = {0,...,N-1}^d.  Corner start \vo = (0,...,0), opposite
%     corner target \va = (N-1,...,N-1) in the corner-to-corner (CC) geometry.
%     (Every formula of the article is stated zero-based and was checked numerically in that
%     form: code/article/closed_form_checks.py, appA_proofs_check.py.)
%   * Time is measured in steps of the discrete-time walk.  "Unit rate" means the
%     continuous-time walk with total jump rate 1 (equivalently q = 1 in q*t).
%   * s WITHOUT subscript is always the Laplace variable; s_k WITH an integer subscript
%     is always sin(pi k / 2N).  The start site is \vo (never "s").
%   * ln is the natural logarithm.
%
% MODEL
%   \dd            d            dimension (1, 2, 3)
%   N              N            number of sites per side
%   \Om            \Omega       lattice {0,...,N-1}^d
%   \nn            n            number of accessible sites (N^d on the clean lattice; |\Cset| with defects)
%   q              q            activity: probability of attempting a move in one step, 0 < q <= 1
%   \vo, \va, \vx  o, a, x      start site, target site, generic site (bold); components o_i, a_i, x_i
%   \Pm            P            transition matrix of the reflecting (cancelled-move) walk; \Pm_1 its q = 1 version
%   \Qm            Q            P with the row and column of the target removed (killed chain)
%   \Lm            L            I - P;   \Lp = L^+  its Moore-Penrose pseudo-inverse
%   \Glap          G            pseudo-inverse of the unit-conductance graph Laplacian (I - P = (q/2d) * Laplacian)
%   \Reff          R            effective resistance (unit resistors);  R_N: opposite corners of the N x N grid
%
% FIRST-PASSAGE QUANTITIES
%   \FPT           \mathcal{T}  random first-passage (hitting) time of the target
%   \Sv(t)         S(t)         survival probability  P(\FPT > t)
%   \pmf(t)        f(t)         first-passage probability mass function  P(\FPT = t) = S(t-1) - S(t)
%   \gd(t)         g(t)         first-passage density of the continuous-time walk (unit rate unless stated)
%   \MF            T            mean first-passage time (MFPT), in steps;  T_{\vo\to\va} when the pair matters
%   \TN            T_N          corner-to-corner MFPT for d = 2;   \TNd{d}  T_N^{(d)} in dimension d
%   \tmode         t^*          mode (most probable first-passage time): smallest argmax_t f(t)
%   \tmodec        t^*_c        mode of the continuous-time density at unit rate (at activity q the mode is t^*_c / q)
%   \tcert         t_cert       time beyond which the PMF is certified to stay below its maximum (Prop. s3_methods:prop-certificate)
%   X_t                         position of the walker after t steps;  h: vector of the mean first-passage times over the start sites
%   \ratio         \varrho      mode-to-mean ratio  t^*/T
%   \Ftil(z)       \widetilde F(z)   probability generating function  sum_t f(t) z^t
%   \Fhat(s)       \widehat F(s)     Laplace transform of g;   \Ptil, \Phat: propagator transforms
%   \hone(m,m')    h_1(m,m')    hitting time from m to m' of the one-dimensional chain with q = 1
%
% SPECTRAL NOTATION
%   \angk{k}       \vartheta_k     = pi k / N
%   \thk{k}{m}     \vartheta_k(m)  = pi k (2m+1) / (2N) = \vartheta_k (m + 1/2)   (zero-based site m)
%   \psi_k(m), \phi_{\mathbf k}(\mathbf x)   one-dimensional cosine eigenvector; product eigenvector of the lattice
%   \ck{k}         c_k          = cos(pi k / 2N)
%   \sk{k}         s_k          = sin(pi k / 2N)
%   \alk{k}        \alpha_k     = 1 (k = 0), 2 (k >= 1)
%   \sigk{k}       \sigma_k     = 2 - cos(\vartheta_k) = 1 + 2 s_k^2
%   \phk{k}        \varphi_k    = arcosh(\sigma_k) = 2 arsinh(s_k)
%   \Bk{k}         \mathcal{B}_k   bracket of form (a) of the single-sum formula
%   \mu_{\mathbf k}             relaxation rates of the reflecting lattice at unit rate, (2/d) sum_i s_{k_i}^2
%   \muone         \mu_1        slowest non-zero one, (2/d) sin^2(pi/2N) ~ pi^2/(2 d N^2)
%   \nu_j, b_j                  decay rates and residues of the first-passage density, g(t) = sum_j b_j exp(-\nu_j t)
%   X, W, A                     X = \mu T (target seen on the diffusive scale), weight and amplitude of the
%                               slowest visible mode (Sec. 6); for CC: W = A = 2d cos^2(pi/2N)
%
% CONSTANTS
%   \gammaE        \gamma_E     Euler's constant
%   \lemn          \varpi       lemniscate constant Gamma(1/4)^2 / (2 sqrt(2 pi)) = 2.6220575542...
%   \Xilem         \Xi          \varpi^4 / \pi^2 = 4.7892679494...
%   C_2, C_0, C_{-2}, ...       coefficients of the large-N expansion of q T_N (d = 2)
%   \Kthree, \Kthreep  K_3, K_3'   q T_N^{(3)} = K_3 N^3 + K_3' N^2 + ...
%   \Ginf          \mathcal{G}_\infty   Green's function of the unit-rate walk on Z^3 (= 6 x the lattice Green's function
%                               of the unit-conductance Laplacian);  \xiH  \xi_H  Hasimoto's constant;  \varkappa (Obs. 4.7)
%   \theta_3(x), \theta_4(x)    Jacobi theta functions with nome argument x:  sum_n x^{n^2},  sum_n (-1)^n x^{n^2}
%   \ell = N - 1/2              half-width of the chain in the continuum limit (Sec. 5.1 only)
%
% DEFECTS (Secs. 7-8)
%   \Bset          \mathcal{B}  set of blocked (inert) sites;   p = |B|/N^2 defect fraction;  M = |B|
%   \Cset          \mathcal{C}  open cluster containing the start (state space of the walk), n = |C|
%   \Thetap        \Theta(p)    time-dilation factor D_0 / D_eff(p);  \Theta_{192}, \Theta_{160}, \Theta_{200}: its three
%                               estimates on tori of that side (Sec. 7.5)
%   \Deff, \Dzero  D_eff, D_0   effective and clean diffusivity
%   \cond, \condz  \Sigma, \Sigma_0   conductivity of the diluted / clean resistor lattice
%   \Pinf          P_\infty     fraction of sites in the accessible cluster
%   \chi_A(\vx)                 single-defect susceptibility of statistic A (A = T or t^*)
%   \perm          \kappa       permeability of an obstacle bond (0 inert, 1 clean)
%   \trap          \Delta       killing probability per step on a partial-trap site
%   \Lgr           \mathcal{L}  unit-conductance graph Laplacian (of the lattice, of Z^2 or of the open cluster)
%   \pk            \mathcal{K}  potential kernel of the planar simple random walk
%   \Pow           \mathcal{W}  dissipated power of a resistor network
%   \Env           \mathcal{E}  envelope of the PMF used to certify the global mode (Sec. 3.4)
%   \Tcap          T_cap        time constant of the slowest geometric term of the PMF ("capture time", Sec. 8)
%   \chiT, \chim   \chi_T, \chi_{t^*}   single-defect susceptibilities of the mean and of the mode
%
% GEOMETRY LABELS (text):  \CC  corner -> opposite corner;  \CtM  corner -> centre target;
%                          \MtC centre -> corner target;    \EXIT centre start, whole boundary layer absorbing
% =====================================================================================

% ---- operators
\DeclareMathOperator{\arcosh}{arcosh}
\DeclareMathOperator{\arsinh}{arsinh}
\DeclareMathOperator*{\argmax}{arg\,max}
\DeclareMathOperator{\Var}{Var}
\newcommand{\Prob}{\mathbb{P}}
\newcommand{\E}{\mathbb{E}}
\newcommand{\dif}{\mathrm{d}}
\newcommand{\eu}{\mathrm{e}}
\newcommand{\iu}{\mathrm{i}}
\newcommand{\T}{\mathsf{T}}                 % transpose
\newcommand{\one}{\mathbf{1}}               % all-ones vector
\newcommand{\Order}{O}

% ---- model
\newcommand{\dd}{d}
\newcommand{\Om}{\Omega}
\newcommand{\nn}{n}
\newcommand{\vo}{\mathbf{o}}
\newcommand{\va}{\mathbf{a}}
\newcommand{\vx}{\mathbf{x}}
\newcommand{\vy}{\mathbf{y}}
\newcommand{\vk}{\mathbf{k}}
\newcommand{\ve}{\mathbf{e}}
\newcommand{\Pm}{P}
\newcommand{\Qm}{Q}
\newcommand{\Lm}{L}
\newcommand{\Lp}{L^{+}}
\newcommand{\Glap}{G}
\newcommand{\Reff}{R}

% ---- first-passage quantities
\newcommand{\FPT}{\mathcal{T}}
\newcommand{\Sv}{S}
\newcommand{\pmf}{f}
\newcommand{\gd}{g}
\newcommand{\MF}{T}
\newcommand{\TN}{T_{N}}
\newcommand{\TNd}[1]{T_{N}^{(#1)}}
\newcommand{\tmode}{t^{*}}
\newcommand{\tmodec}{t^{*}_{\mathrm{c}}}
\newcommand{\tcert}{t_{\mathrm{cert}}}
\newcommand{\ratio}{\varrho}
\newcommand{\Ftil}{\widetilde{F}}
\newcommand{\Fhat}{\widehat{F}}
\newcommand{\Ptil}{\widetilde{P}}
\newcommand{\Phat}{\widehat{P}}
\newcommand{\hone}{h_{1}}

% ---- spectral notation
\newcommand{\angk}[1]{\vartheta_{#1}}
\newcommand{\ck}[1]{c_{#1}}
\newcommand{\sk}[1]{s_{#1}}
\newcommand{\alk}[1]{\alpha_{#1}}
\newcommand{\thk}[2]{\vartheta_{#1}(#2)}
\newcommand{\sigk}[1]{\sigma_{#1}}
\newcommand{\phk}[1]{\varphi_{#1}}
\newcommand{\Bk}[1]{\mathcal{B}_{#1}}
\newcommand{\muone}{\mu_{1}}

% ---- constants
\newcommand{\gammaE}{\gamma_{\mathrm{E}}}
\newcommand{\lemn}{\varpi}
\newcommand{\Xilem}{\Xi}
\newcommand{\Kthree}{K_{3}}
\newcommand{\Kthreep}{K_{3}'}
\newcommand{\Ginf}{\mathcal{G}_{\infty}}
\newcommand{\xiH}{\xi_{\mathrm{H}}}

% ---- defects
\newcommand{\Bset}{\mathcal{B}}
\newcommand{\Cset}{\mathcal{C}}
\newcommand{\Thetap}{\Theta}
\newcommand{\Deff}{D_{\mathrm{eff}}}
\newcommand{\Dzero}{D_{0}}
\newcommand{\cond}{\Sigma}
\newcommand{\condz}{\Sigma_{0}}
\newcommand{\Pinf}{P_{\infty}}
\newcommand{\perm}{\kappa}
\newcommand{\trap}{\Delta}
\newcommand{\Lgr}{\mathcal{L}}
\newcommand{\pk}{\mathcal{K}}
\newcommand{\Pow}{\mathcal{W}}
\newcommand{\Env}{\mathcal{E}}
\newcommand{\Tcap}{T_{\mathrm{cap}}}
\newcommand{\chiT}{\chi_{\MF}}
\newcommand{\chim}{\chi_{\tmode}}

% ---- geometry labels (text mode)
\newcommand{\CC}{\textsc{cc}}
\newcommand{\CtM}{\textsc{c2m}}
\newcommand{\MtC}{\textsc{m2c}}
\newcommand{\EXIT}{\textsc{exit}}

% ---- theorem environments (one counter per section; appendices: A.1, B.1, ...)
\theoremstyle{plain}
\newtheorem{theorem}{Theorem}[section]
\newtheorem{proposition}[theorem]{Proposition}
\newtheorem{lemma}[theorem]{Lemma}
\newtheorem{corollary}[theorem]{Corollary}
\newtheorem{conjecture}[theorem]{Conjecture}
\theoremstyle{definition}
\newtheorem{definition}[theorem]{Definition}
\newtheorem{observation}[theorem]{Numerical observation}
\newtheorem{heuristic}[theorem]{Formal result}
\theoremstyle{remark}
\newtheorem{remark}[theorem]{Remark}
% computer-assisted statements (Sections 4 and 7): printed as "Theorem 7.n (<title>; computer-assisted)"
\theoremstyle{plain}
\newtheorem{catheoreminner}[theorem]{Theorem}
\newtheorem{cacorollaryinner}[theorem]{Corollary}
\newenvironment{catheorem}[1]{\begin{catheoreminner}[#1; computer-assisted]}{\end{catheoreminner}}
\newenvironment{cacorollary}[1]{\begin{cacorollaryinner}[#1; computer-assisted]}{\end{cacorollaryinner}}
\theoremstyle{definition}
\newtheorem{refuted}[theorem]{Refuted statement}
% pointers into the companion notes R0-R5 (folder proofs/ of the repository) and status markers of Section 7
\newcommand{\Sref}[2]{[\textsf{#1},~#2]}
\newcommand{\lateproof}{\ensuremath{^{\star}}}     % proof written after the two internal checks of its companion note (read line by line afterwards)
\newcommand{\combined}{\ensuremath{^{\diamond}}}   % deduction combining two companion notes, made in R0
\newcommand{\tcf}{t_{\mathrm{cf}}}                 % closed form T ln(2dX)/(X+2d-1), Eq. (s6_mechanism:eq-L1-simple)

% ---- status vocabulary used in the text (keep the wording uniform)
%   Theorem / Proposition / Lemma / Corollary : complete proof in the article.
%   Theorem (computer-assisted) (env. catheorem, cacorollary): proof reduced by hand to finitely many inequalities
%                                               verified in exact or ball arithmetic; trusted base named (Sec. 7.1).
%   Refuted statement                         : a natural conjecture (suggested by d = 1 or by numerics) that is false.
%   Formal result (env. heuristic)            : derived by formal asymptotics, verified numerically, not proved.
%   "closed-form approximation"               : ALWAYS Eq. (s6_mechanism:eq-L1) (its simplified form: eq-L1-simple);
%   "exact two-/three-exponential truncation" : maximiser of the two or three slowest exponentials with numerically
%                                               exact rates and residues (a different approximation; never "two-pole").
%   Numerical observation                     : established by exact computation at the sizes stated only.
%   Conjecture                                : believed, supported by evidence, unproved.
